\section{Introduction}
\label{sec:intro}

Linear probes read a language model's internal activations directly, which makes them attractive
monitors for deception: a model that lies need not say anything that gives it away
\citep{goldowskydill2025detecting}. But a probe is trained on one distribution and deployed on
others, and probes often fail to transfer
\citep{kirch2025offpolicy,natarajan2026targeted,kumar2026pressure}. Liars' Bench makes the failure
concrete \citep{kretschmar2025liars}. Across seven datasets of lies generated on-policy by four
open-weight models, a probe trained on instructed honest and deceptive responses, following
\citet{goldowskydill2025detecting}, reaches balanced accuracy close to chance on most datasets. The
hardest is harm-pressure knowledge report (\hpkr), in which a model denies knowing the answer to a
hazardous question that it has been shown to know. Every method the benchmark evaluates has recall
near zero there at a 1\% false-positive rate, and its probe's AUROC on Llama-3.3-70B is
\V{lb-mean-auroc-hpkr}.

\citet{mahadik2026persona} propose a different source of directions. Following the Assistant Axis
methodology \citep{lu2026assistant}, they prompt a model to play deceptive roles, such as a spin
doctor or a coverup operator, and honest ones, such as an ombudsman or a fact checker. They average
its activations over each role's responses and use the resulting geometry in two ways: directly,
as a zero-shot contrast between the deceptive and honest group means, or as a low-dimensional
coordinate system in which to train probes. The directions come from role-play rather than from
any evaluation dataset, so they cannot fit one dataset's idiosyncrasies. On Llama~3.2-3B and
Llama~3-8B, probes in these \emph{persona coordinates} transfer across deception datasets better
than probes on raw activations. Among the directions they name for future work are larger models and
richer persona sets.

We take both. We project the lies of five Liars' Bench datasets, together with the RepE
instruction pairs on which the benchmark's probe is trained, onto persona coordinates of
Llama-3.3-70B-Instruct at \V{depth-n} depths, and measure cross-dataset transfer over the full
train--test matrix (\cref{fig:pipeline}). This exposes a failure the original study did not test.
On \hpkr, the published persona contrast inverts: it ranks lies as more honest than honest answers,
with AUROC \V{zs-dec-hpkr-L20} at layer~20. The cause lies in the personas. \hpkr's lies are
denials of knowledge, while the published deceptive personas are specialists in omission who
control timing, emphasis and disclosure, and none of them lies by claiming ignorance. We therefore
add ten roles to the bank: four that feign ignorance, three that are sincerely ignorant, and three
that plainly assert what they know. Projected into these \emph{knowledge-denier persona
coordinates}, the lies separate.

\begin{figure}[t]
\centering
% The pipeline schematic of Figure 1.
\begin{tikzpicture}[
    node distance=0.42cm and 0.30cm,
    box/.style={draw=black!40, rounded corners=2pt, align=center, font=\footnotesize,
                inner sep=3pt, minimum height=1.15cm, text width=2.78cm, fill=#1},
    box/.default=white,
    arrow/.style={-{Stealth[length=4.5pt]}, draw=black!60, semithick}]
  \node[box=orange!12] (personas) {\textbf{36 personas}\\5 prompt variants\\30 questions each};
  \node[box, right=of personas] (gen) {Llama-3.3-70B\\role-plays each\\(\V{bank-generated} rollouts)};
  \node[box, right=of gen] (judge) {Judge keeps the\\in-character\\rollouts (\V{bank-kept})};
  \node[box, right=of judge] (vec) {Persona vector:\\mean activation\\over its responses};
  \node[box=blue!9, right=of vec] (coords) {\textbf{Persona}\\\textbf{coordinates}\\$d$ and leading PCs};

  \node[box=black!5, below=0.75cm of personas] (example) {\textbf{Liars' Bench}\\response (or RepE,\\Alpaca control)};
  \node[box, right=of example] (acts) {Llama-3.3-70B\\activations at\\layer 20, 40 or 44};
  \node[box, right=of acts] (pool) {Mean over the\\response tokens:\\$z$};
  \node[box, right=of pool] (score) {Project onto $d$:\\[2pt]$s(z)=\langle z,\hat d\,\rangle$};
  \node[box=blue!9, right=of score] (eval) {\textbf{Lie score}\\zero-shot, or as a\\probe feature};

  \foreach \a/\b in {personas/gen, gen/judge, judge/vec, vec/coords,
                     example/acts, acts/pool, pool/score, score/eval}
    \draw[arrow] (\a) -- (\b);
  \draw[arrow] (coords.south) |- ($(score.north)+(0,0.22)$) -- (score.north);
\end{tikzpicture}

\caption{Projecting lies onto persona coordinates. \textbf{Top:} the model role-plays each persona
under five prompt variants and 30 shared questions. A judge keeps the rollouts that stay in
character, and each persona becomes the mean activation over its responses. The coordinates, a
contrast direction $d$ between deceptive and honest personas and the leading principal components,
never see an evaluation dataset. \textbf{Bottom:} each Liars' Bench response is mean-pooled and
projected onto $d$. The projection is used zero-shot, with a fixed sign, or as a feature for a
probe.}
\label{fig:pipeline}
\end{figure}

\paragraph{Contributions.} In this work, we make the following contributions:
\begin{itemize}[leftmargin=1.2em, itemsep=2pt, topsep=3pt]
\item \textbf{A persona bank for deception by denial.} We extend the 26 personas of
  \citet{mahadik2026persona} with \V{bank-new} roles: four that feign ignorance, three that are
  sincerely ignorant, and three that plainly assert what they know. From them we derive a
  knowledge-denial contrast (4 deceptive vs.\ 6 honest personas) and an extended contrast that
  pools both deceptive and both honest sides (12 vs.\ 13).
\item \textbf{Persona coordinates at 70B.} We scale the method from 3B and 8B models to
  Llama-3.3-70B-Instruct, evaluated on five Liars' Bench datasets and RepE at \V{depth-n} layers,
  every fourth from \V{depth-first} to 76 plus the last block.
\item \textbf{A better score on Liars' Bench's hardest dataset.} On \hpkr, with no benchmark
  labels, the knowledge-denial contrast reaches \V{zs-den-hpkr-L40} AUROC and
  \V{op-den-balacc-hpkr-L40} balanced accuracy at 1\% FPR at layer~40. For the same model, Liars'
  Bench reports \V{lb-mean-balacc-hpkr} for its Mean Probe, \V{lb-followup-balacc-hpkr} for its
  Follow-Up Probe, and \V{lb-upper-balacc-hpkr} for a probe trained on privileged in-distribution
  labels. A RepE-trained probe in our own pipeline stays at \V{bp-balacc-hpkr-L40} at every one of the
  \V{depth-n} layers, though from layer~\V{depth-hpkr-base-auroc-peak-layer} it ranks \hpkr{}
  better than any contrast we build (\cref{sec:depth}).
  We also identify why the published contrast inverts there: none of its deceptive personas lies
  by denying knowledge.
\item \textbf{Cross-dataset transfer at 70B.} Across the full train--test matrix, the zero-shot
  extended contrast reaches \V{off-ext-L40} mean off-diagonal AUROC at layer~40, against
  \V{off-base-L40} for 8192-dimensional probes each trained on one dataset, and leads them at every
  one of the \V{depth-n} layers we swept, by at least \V{depth-off-ext-over-base-min} AUROC. Those
  probes lose \V{depth-base-drop-min}--\V{depth-base-drop-max} AUROC on leaving the dataset they
  were fitted to; the contrast, fitted on none, loses at most \V{depth-ext-drop-max}. It beats their average
  over training sets on every held-out dataset at every layer through
  \V{depth-pertest-holds-last}; deeper, where every contrast is weakest,
  \V{depth-pertest-exception} are exceptions (\cref{sec:depth}).
\item \textbf{One direction beats the persona subspace.} The 1-d contrast outscores probes on the
  three leading persona principal components at \V{depth-pc-wins-complement} of the \V{depth-n}
  layers we swept, by \V{gap-ext-extpcthree-L20} AUROC at layer~20 and
  \V{gap-ext-extpcthree-L44} at layer~44; it reverses only inside the deep trough
  (\cref{sec:depth}).
\item \textbf{An evaluation under the benchmark's own protocol.} We compare the persona contrasts
  with the benchmark's RepE-trained probe at its 1\% false-positive threshold on Alpaca, with
  bootstrap intervals. The appendices report how the direction behaves at the Follow-Up Probe's
  token and when read from individual response tokens.
\item \textbf{A verified, reproducible pipeline.} Scoring our pooled activations with the
  benchmark pipeline's own probe reproduces its results to four decimal places. Our measurement of
  the Follow-Up Probe, \V{fu-measured-auroc} mean AUROC over the five datasets, matches the
  published \V{lb-followup-auroc-avg}. We release the code and results.%
  \footnote{\url{https://github.com/zami456/expLB}}
\end{itemize}

Two caveats frame these results, and \cref{sec:limitations} develops both. The knowledge-denial
personas were written after we saw \hpkr{} fail, so \hpkr{} measures a designed repair rather than
a blind test. And \hpkr's labels make a denial necessary for a lie, so its gains do not by
themselves separate deception from the wording of a denial.
